\chapter{Divide and Conquer}
\label{chap:divide-and-conquer}

% Closed question: How can an instrument compare records and repeat the result
% without pretending the comparison is the thing measured?

\begin{chapteressay}

A speedometer reads sixty miles per hour. The driver glances at the
dial, confirms the reading, and looks back at the road. Three seconds
later the dial still reads sixty. The two readings agree. The driver
does not interpret the agreement as evidence that the speedometer's
reading is the speed. The agreement is evidence that the speedometer
is repeating its comparison reliably.

The speed is something else. The speed is what the car is doing. The
speedometer's reading is the algorithm's report on its comparison
between the rotation rate of the drive shaft and the calibrated
reference. The reading and the speed are not the same object. The
reading is the output of the comparison; the speed is the physical
event the comparison was designed to characterize.

This distinction is the chapter's subject. An instrument compares.
The comparison's output is the reading. The reading is not the
measurement; the measurement is what the reading is interpreted as
indicating. The interpretation depends on calibration, context, and
the chain of comparisons that produced the reading. To say the
speedometer measures the speed is shorthand for: the speedometer's
internal mechanism performs a sequence of admissible comparisons, the
sequence terminates with a final comparison whose output is the
displayed number, and the displayed number is interpreted as the
speed by virtue of the speedometer's calibration. The reading is the
algorithm's output; the speed is the algorithm's intended referent.

When the algorithm is divide-and-conquer, the structure becomes
explicit. A merge sort on a million records does not measure any
individual record; it produces a permutation. A binary search on a
sorted array does not measure the target; it produces the target's
index. An FFT does not measure the signal; it transforms the
signal's representation. A tournament sort does not measure the
runners; it identifies a winner. In each case, the algorithm's
output is the produced structure (the permutation, the index, the
spectrum, the winner). The comparisons inside the algorithm are how
the structure is produced; they are not the structure itself.

This is the algorithmic statement of an old metrological principle.
The instrument and the measurement are distinct. The instrument
performs a procedure; the procedure produces a reading; the reading
is interpreted as a measurement. The procedure is the algorithm.
Different procedures (different algorithms) on the same input
produce different readings. Different calibrations of the same
procedure produce different interpretations of the same reading.
The structure of the algorithm and the structure of the calibration
together determine what the reading means.

The chapter develops this through five algorithmic forms, each one a
different way of comparing records while keeping the comparison
distinct from the measurement.

The five forms are not arbitrary. They are the canonical
divide-and-conquer algorithms: binary search (locate one element in
a sorted structure), merge sort (impose order on an unordered
structure), FFT (transform between representations), tournament sort
(extract a distinguished element), and the typeclass machinery
\texttt{REPEATABLE} that certifies all of the above. The chapter's
claim is that this family is forced by the requirement of
repeatability. Any repeatable comparison structure factors through
some combination of these primitives.

The second concrete example, beyond the speedometer, is the mail
sorting machine. A mail-processing center handles a million envelopes
per day. Each envelope's address is read by an optical scanner. The
scanner's reading is a multi-digit zip code. The center then sorts
envelopes by zip code into delivery sequences. The comparisons inside
the sort are not the address; the address is input. The output is
the delivery sequence: the permutation of envelopes that the carrier
will deliver in order.

A million envelopes produce roughly ten million comparisons. The
sorting machine produces about one hundred forty delivery sequences.
The comparison-to-output ratio is about seventy thousand to one. None
of the comparisons is a delivery sequence; the delivery sequences are
the outputs the comparisons produce.

If the machine's comparisons are not \texttt{REPEATABLE} --- if the
same envelope's zip code is sometimes read as 90210 and sometimes as
90211 --- the delivery sequences will be wrong. The machine's
correctness depends on the underlying scanner being repeatable. A
miscalibrated scanner produces misrouted mail.

The chapter's closed question is:

\begin{center}
\emph{How can an instrument compare records and repeat the result
without pretending the comparison is the thing measured?}
\end{center}

The answer the chapter develops is divide-and-conquer with
\texttt{REPEATABLE} comparisons. The instrument is the algorithm.
The comparisons are the algorithm's primitive operations. The
\texttt{REPEATABLE} typeclass certifies that the comparisons
produce consistent outputs. The algorithm composes the comparisons
into a structured output (a permutation, an index, a spectrum, a
winner). The structured output is the measurement. The comparisons
are not the measurement; the comparisons are how the measurement
is produced.

This chapter is short by design. The argument is straightforward
once the algorithmic form is named. Subsequent chapters build on it:
Chapter 3 examines what the instrument carries between comparisons
(the metavariable, the unresolved state). Chapter 4 examines how
multiple records combine (union-find). Chapter 5 examines distance
as count of refinements (bisection and the spline). Chapter 6 is
the load-bearing chapter, examining the boundary where the algorithm
runs out: the selection rule requires exact Kolmogorov information,
and exact $K$ is not computable. The earlier chapters build the
algorithmic vocabulary the noncomputability argument requires.

\end{chapteressay}

\section{Binary Search}
\label{se:binary-search}

A sorted array of one million integers can be searched in twenty comparisons.
Sequential scan would require, on average, half a million comparisons. The
difference is the algorithm. Binary search exploits the sorting to eliminate
half the remaining candidates with each comparison.

The procedure is simple. Maintain two indices, $\text{low}$ and $\text{high}$,
initialized to the first and last positions of the array. While
$\text{low} \le \text{high}$, compute the middle index $m = \lfloor
(\text{low} + \text{high}) / 2 \rfloor$. Compare the array element at
position $m$ to the target. If the element equals the target, return $m$.
If the element is less than the target, set $\text{low} = m + 1$ (the
target, if present, is in the upper half). If the element is greater than
the target, set $\text{high} = m - 1$ (the target is in the lower half).
The loop terminates either when a match is found or when
$\text{low} > \text{high}$, indicating the target is not in the array.

Each iteration of the loop halves the size of the remaining search space.
Starting with $n$ elements, after $k$ iterations the search space has size
$n / 2^k$. The loop terminates when the search space has size zero or
one. The number of iterations is therefore at most $\lceil \log_2 n
\rceil$. For $n = 10^6$, this is twenty iterations.

The algorithm's correctness depends on two preconditions. First, the array
must be sorted. Without sorting, binary search produces incorrect results
because the comparison at position $m$ does not constrain the contents of
the upper and lower halves in the required way. Second, the comparison
must be a total order: every pair of elements must be comparable as less,
equal, or greater. Both preconditions are part of the instrument's
calibration. The sorting is the precomputation; the total order is the
typeclass constraint.

The output of binary search is the index of the target if present, or the
not-found indicator. It is not the target itself; the target was already
known when the search began. The output is the location in the record
where the target lives. This is the operational distinction at the heart
of the chapter's question: the comparison is not the thing measured; the
comparison is the procedure for locating it.

What is measured? The location. The target is the query; the array is the
record; the location is the answer. The comparison is the algorithm that
maps query and record to answer. Saying ``the target equals the element at
position $m$'' is not a measurement of the target. The target's value was
input data, not an output. What was produced was the index $m$, which is
information about the array's structure, not about the target.

This distinction is crucial when the comparison is being performed for the
first time on a particular pairing. A scientist using a spectrometer to
identify an unknown compound is not measuring the compound's spectrum --- the
spectrum is the instrument's output, the raw data. The measurement is the
comparison of that spectrum against a library of reference spectra,
yielding an identification (the index in the reference library of the
closest match). The spectrometer produces the spectrum; the database
lookup produces the identification. The instrument is the composition.

In Lean, binary search appears as the typeclass \texttt{BINARY} in
\texttt{Episode2.lean}. The class names what the comparison produces: a
binary output (one of three values: less, equal, greater). The comparison
itself is a function from a pair of values to this three-valued type.
\texttt{Trial.le} and \texttt{Trial.lt} are the typeclasses for less-than
and less-than-or-equal trials. Together they form a decidable total order.

The reason binary search is the chapter's first algorithm is that it is
the simplest divide-and-conquer procedure. The recursion is trivial: each
call halves the problem. The comparison at the midpoint is the entire
work of the recursive step. The base case is when the size drops to zero
or one. The recursion is also tail-recursive: the same loop can be
implemented iteratively without a function call stack.

Iterative binary search is also the cheapest divide-and-conquer
algorithm. The constants in the $O(\log n)$ bound are small. A single
comparison plus a single integer divide plus two integer updates per
iteration. Twenty iterations for $n = 10^6$. Memory usage is $O(1)$: no
additional storage beyond the indices and the comparison variable. This
makes binary search a primitive in many higher-level algorithms: it is
the lookup operation in B-trees, in skip lists, in sorted associative
arrays. Whenever an ordered collection needs a fast lookup, binary
search is the default mechanism.

The compiler's analog is namespace lookup with a sorted dictionary. When
the elaborator sees the name \texttt{Nat.add}, it walks the namespace
tree. Within each node, the children are stored in sorted order. The
lookup at each level is binary search on the children's names. The cost
of resolving a fully qualified name with $k$ dotted components, in a
namespace with $n$ entries at each level, is $O(k \log n)$. For Mathlib,
which has hundreds of thousands of names, this is the difference between
sub-millisecond name resolution and seconds-long delays. The Lean
compiler's responsiveness depends on this algorithmic choice.

Binary search is also the algorithmic form of a decision procedure with
gradual information gain. Each comparison provides one bit of information
about the target's location (it is in the upper half or the lower half).
After $k$ comparisons, the search has narrowed the candidates from $n$ to
$n / 2^k$. The information content of the answer is $\log_2 n$ bits ---
exactly the index of the matched position, interpreted as a binary
integer. The algorithm extracts these bits one at a time. This is the
ideal case: the information content of the output matches the work done
to produce it. Binary search is information-theoretically optimal for
sorted-array search.

\section{Merge Sort}
\label{se:merge-sort}

Binary search assumes a sorted array. The next question is how the array
became sorted. The canonical algorithm is merge sort.

Merge sort sorts a list of $n$ elements in three steps. First, if $n \le 1$,
the list is already sorted; return it. Second, split the list into two
halves, of size $\lceil n/2 \rceil$ and $\lfloor n/2 \rfloor$. Recursively
sort each half. Third, merge the two sorted halves into a single sorted
list. The merge step compares the front elements of the two halves,
selects the smaller, and emits it; it then advances the corresponding
half's index and repeats until both halves are exhausted.

The recursion is the divide. The merge is the conquer. The recursion depth
is $\log_2 n$. The merge step at each recursion level processes a total of
$n$ elements (because the two halves at level $k$ together have $n$
elements between them). The total work is $O(n \log n)$, which is
asymptotically optimal for comparison-based sorting.

The algorithm's structure displays the chapter's question precisely.
Sorting a list of one million records by surname is not a measurement of
any individual record. The output is a permutation: a reordering of the
input that satisfies the order relation. The comparisons made during the
sort produce no record of any individual surname; they produce the
permutation that puts the surnames in the right sequence. The
permutation is the measured output. Each comparison is a step in the
algorithm; the algorithm is the instrument; the permutation is the
reading.

Each comparison in merge sort is a Boolean trial: is the front of the
left half less than or equal to the front of the right half? The answer
selects which element is emitted next. The trial is the
\texttt{BINARY} typeclass. The result is one of two outputs: emit-left
or emit-right. The trial itself does not produce information about the
absolute position of either element in the final permutation; it
produces only the local information that one element comes before
another. The local information accumulates, over the course of the
merge, into the global information of the sorted order.

This is a precise instance of what the chapter is naming. The
comparison is the algorithm, not the measurement. The measurement is the
sorted output. The repeatability of the comparison --- the same two
inputs always producing the same Boolean output --- is what makes the
sorted output deterministic. If the comparison were nondeterministic,
running the sort twice on the same input could produce different
permutations.

Merge sort is also a stable sort. If two elements compare equal, their
relative order in the output matches their relative order in the input.
Stability is a property of the algorithm, not of the comparison. The
comparison says only ``these two elements are equal''; the algorithm
chooses, when two equal elements are encountered, to emit the one from
the left half before the one from the right half. This choice is the
algorithm's commitment to preserving prior order. The algorithm is the
instrument; the comparison is the trial; the stability is the
instrument's calibration choice.

The cost in memory is $O(n)$ for the merge buffers. This is the standard
implementation; in-place merge sort exists but is more complex and
slower in practice. The memory cost reflects the fact that the merge
step cannot operate purely on the input; it requires temporary storage
to hold the half-sorted intermediate results. The instrument needs
scratch space.

The Lean analog is \texttt{Study} in \texttt{Episode2.lean}. A
\texttt{Study} is a structured inquiry: a sequence of trials that
together establish a comparison conclusion. \texttt{Study.le} is the
partial order on studies: one study is at most another when its
conclusion is contained in the other's. A merge sort, viewed as a
study, has $n \log n$ trials and produces, as its conclusion, the
permutation. Each individual trial is a \texttt{Trial.le} instance; the
study composes them into the larger result.

\texttt{REPEATABLE} in \texttt{Episode2.lean} formalizes the property
that a sort produces the same permutation when run twice on the same
input. The typeclass requires that the underlying comparison be
deterministic: \texttt{Trial.le a b} is a function of $a$ and $b$ only,
not of any other state. If the comparison consults a clock, a random
number generator, or any external resource, the typeclass instance
cannot be provided, and the sort cannot be certified as repeatable.

The compiler analog is the elaboration of a definition. When Lean
elaborates a complicated definition with many implicit arguments, it
performs a sequence of unification steps, each one a small comparison
between expected and actual types. The order of unification is
determined by the elaboration strategy. Two runs of the elaborator on
the same input produce the same result --- the same set of unification
decisions, in the same order --- because each unification step is
deterministic. The elaborator is a sort: it produces a structured
output (the elaborated term) by repeatedly applying a deterministic
comparison (unification) to an input set of constraints. The
comparison is not the elaboration; the elaboration is the algorithm
that uses the comparison.

A practical observation: merge sort scales well to external sorting.
When the input does not fit in main memory, the algorithm can be run
on disk-backed files, with each half sorted as a separate file and the
merge step streaming the two files into a third. This is how large
database systems sort tables that exceed available RAM. The
algorithm's divide-and-conquer structure makes external sorting
feasible: the recursion can be unrolled into a sequence of merge
passes over the data, with each pass reading and writing the entire
dataset once. The cost is $O(n \log_k n)$ I/O operations for a
$k$-way merge, which for typical disk speeds is acceptable on
terabyte-scale data.

The classical achievement of merge sort is its provable $O(n \log n)$
lower bound. Any comparison-based sort that always produces the
correct output must, in the worst case, make $\Omega(n \log n)$
comparisons. The lower bound is information-theoretic: the number of
distinct permutations of $n$ elements is $n!$, and distinguishing
among $n!$ outcomes requires $\log_2 (n!) \approx n \log_2 n$ bits of
information. Each comparison provides at most one bit. Therefore the
number of comparisons is at least $n \log_2 n$. Merge sort achieves
this bound up to a constant. The algorithm is, in this precise sense,
optimal.

\section{FFT: Recursive Divide}
\label{se:fft-recursive-divide}

The Discrete Fourier Transform (DFT) of a length-$N$ signal $x_0, x_1,
\ldots, x_{N-1}$ produces a length-$N$ frequency spectrum $X_0, X_1, \ldots,
X_{N-1}$, computed by
\[
X_k = \sum_{n=0}^{N-1} x_n \, e^{-2\pi i \, k n / N}.
\]
The naive computation of each $X_k$ requires $N$ multiplications and $N-1$
additions; computing all $N$ values costs $O(N^2)$ operations. For
$N = 10^6$, this is $10^{12}$ operations, which is hours on a typical CPU.

The Fast Fourier Transform, due to Cooley and Tukey in 1965 (though
discovered earlier by Gauss and others and forgotten), reduces this to
$O(N \log N)$ operations. For $N = 10^6$, the cost is approximately
$2 \times 10^7$ operations: a fraction of a second. The speedup is six
orders of magnitude, achieved by a divide-and-conquer factoring of the
sum.

The factoring exploits an algebraic identity. When $N$ is a power of two,
the sum can be split into even-indexed and odd-indexed terms:
\[
X_k = \sum_{n=0}^{N/2-1} x_{2n} \, e^{-2\pi i \, k (2n) / N}
     + \sum_{n=0}^{N/2-1} x_{2n+1} \, e^{-2\pi i \, k (2n+1) / N}.
\]
The first sum is a DFT of length $N/2$ on the even-indexed samples. The
second sum is a DFT of length $N/2$ on the odd-indexed samples, multiplied
by a twiddle factor $e^{-2\pi i \, k / N}$. The full transform is the
combination of two half-length transforms plus $N$ complex multiplications
and additions (the butterfly stage).

The recurrence is $T(N) = 2 T(N/2) + O(N)$. By the master theorem, this
solves to $T(N) = O(N \log N)$. The depth of recursion is $\log_2 N$.
Each level does $O(N)$ work in the butterfly combining step. For
$N = 2^{20} \approx 10^6$, the recursion has twenty levels, and each
level does about $10^6$ operations. The total is roughly $2 \times 10^7$,
matching the asymptotic estimate.

The FFT is divide-and-conquer at maximum information density. Each
butterfly operation is a $2 \times 2$ matrix multiplication: two complex
inputs produce two complex outputs. The matrix is
\[
\begin{pmatrix}
1 & \omega \\
1 & -\omega
\end{pmatrix},
\]
where $\omega = e^{-2\pi i k / N}$ is the twiddle factor specific to the
position in the recursion. The butterfly is the smallest unit of FFT
computation: two inputs, two outputs, one twiddle factor, four complex
multiplications, two complex additions. The entire FFT is $N \log N / 2$
butterflies.

Each butterfly is an admissible trial: it produces a comparison between
two frequency components, separating their sum and difference. The trial
is deterministic: the same two inputs and the same twiddle factor produce
the same two outputs. Repeated application of the FFT to the same signal
produces the same spectrum every time. This is the chapter's
\texttt{REPEATABLE} property at the level of the largest commonly-used
algorithm in numerical computing.

What is measured by the FFT? The frequency content of the signal. The
output $X_k$ is the amplitude (complex-valued) of the $k$-th harmonic.
Each $X_k$ is a record of the signal's projection onto a specific
frequency. The FFT does not change the signal; it transforms the
representation. The same physical signal, expressed in time domain
($x_n$) or frequency domain ($X_k$), is the same physical event. The FFT
is the algorithmic instrument that converts between the two
representations.

This connects FFT to the chapter's question. The transformation is not
the measurement of the signal; the signal was already known. The
transformation is a structural reorganization that exposes a different
aspect of the same record. A signal that looks random in the time domain
may reveal periodic components in the frequency domain. The FFT is the
instrument that makes those components visible. The measurement is what
the instrument shows; the algorithm is how the instrument processes the
input.

The FFT is also the canonical example of an algorithm that produces a
unique minimal-information representation. The DFT is invertible: from
$X_k$, the original $x_n$ can be recovered by an inverse FFT of the
same complexity. The transform preserves all information. It is a
unitary change of basis. The frequency domain representation is not
``richer'' or ``poorer'' than the time domain representation; they
contain the same information, expressed differently. The minimum-
information property is what makes the FFT useful: it lets the user
work in whichever representation is most convenient for the task.

In Lean, the FFT analog is the typeclass cascade that handles
elaboration. Every implicit argument unification is, in effect, a small
transform: the elaborator reorganizes the expression to expose the
constraints that need to be solved. The cascade of unifications has the
same divide-and-conquer structure as the FFT: at each level, the
elaborator splits the constraint into smaller constraints, solves them
recursively, and combines the results. The compiler's elaboration cost,
when expressed in heartbeats, is approximately $O(n \log n)$ in the size
of the term being elaborated. The same algorithmic form, the same
asymptotic cost.

A practical observation: the FFT's $O(N \log N)$ cost made many
applications feasible that were previously not. Real-time spectrum
analysis on audio signals, image compression in JPEG, fast convolution
of large signals, polynomial multiplication in $O(N \log N)$ rather than
$O(N^2)$, and many others. The algorithmic improvement transformed entire
fields of engineering. The chapter's point is not that the FFT is fast;
it is that the FFT's structure is the algorithmic embodiment of the
chapter's claim. Divide and conquer is not an optimization heuristic; it
is the structural form forced by the chapter's question about repeatable
comparison.

\section{Tournament Sort}
\label{se:tournament-sort}

A field of sixty-four runners is to produce one winner. Each pairwise race is
a definite outcome: one runner finishes ahead of the other. The simplest
fair procedure is a single-elimination tournament. Round one: thirty-two
races; the winners advance, the losers are eliminated. Round two: sixteen
races. Round three: eight. Round four: four. Round five: two. Round six:
one. After six rounds and sixty-three races, one runner remains: the
champion.

Tournament sort generalizes this. Given $n$ elements and a comparison
operator, build a complete binary tree with $n$ leaves. Place the elements
at the leaves. At each internal node, store the winner (the larger or
smaller, depending on the desired order) of its two children. Propagate
upward: the root contains the overall winner. This takes $n - 1$
comparisons.

To extract the next element after the winner, mark the winner's leaf as
removed and re-run the comparisons along the path from that leaf to the
root. The path has $\log_2 n$ edges, so $\log_2 n$ comparisons restore the
tree. The next winner is at the root. Repeat to extract all elements in
sorted order. Total comparisons: $n - 1$ for the initial tournament plus
$(n-1) \log_2 n$ for the extractions. The total is $O(n \log n)$,
asymptotically the same as merge sort.

Tournament sort's value is not its asymptotic cost --- merge sort matches
it. The value is that the algorithm finds the winner without producing the
full sorted order. Sixty-three races identify the fastest runner; the
ranking of the second-place finisher requires an additional six races
(re-running the path through the second-place runner's bracket). The
algorithm produces ranked information incrementally, on demand, rather
than computing the full sorted order eagerly.

This matters for the chapter's question. Tournament sort is the
algorithmic form of a comparison whose result is not the full ranking but
a specific named question: who is the winner? The comparison is performed
exactly enough times to answer the question, and no more. The remaining
elements are partially ordered (we know they all lost to someone in the
bracket) but not fully ordered (we do not know their relative ranks).

The comparison is the algorithm. The winner is the measurement. The
algorithm uses sixty-three comparisons; the measurement is one runner's
identity. The comparisons are not the measurement. The measurement is the
question the tournament was designed to answer.

What does this teach? It teaches that the comparison's purpose is
incidental to the comparison itself. Sixty-three runners participate in
the tournament, each of them comparing against another runner in some
round. Each runner's individual experience is a comparison. The aggregate
experience produces the winner. The aggregate is the measurement; the
individual comparisons are the procedure.

If the tournament asks ``who is the fastest runner?'', the answer is one
name. If the tournament asks ``what is the full ranking?'', the answer is
sixty-four names in order. The two questions require different
algorithms. The tournament-as-described answers the first question in
$O(n)$ comparisons. Answering the second requires $O(n \log n)$ further
comparisons. The questions are different; the algorithm is different; the
measurements are different. The comparisons in each case are not the
measurement; they are how the algorithm progresses toward producing the
measurement.

Tournament sort also exposes a subtle property: the comparisons in
different parts of the tree are independent and can be parallelized. The
thirty-two first-round races can all be run simultaneously --- on
thirty-two separate tracks, or on thirty-two CPU cores. The sixteen
second-round races can follow, but they cannot begin until the first
round completes. The total wall-clock time is $\log_2 n$ rounds, not $n -
1$ sequential races. For $n = 64$, the wall-clock time drops from
sixty-three rounds to six rounds. The algorithm exploits parallelism
naturally.

This is why tournament sort, and divide-and-conquer in general, scales
well on parallel hardware. The recursion's branching gives independent
subproblems at each level. The combining stage (which in tournament sort
is the comparison at the parent node) is the only sequential dependency.
For algorithms with cheap combining steps (like tournament sort), the
parallel speedup is nearly linear in the number of processors.

The Lean analog of tournament sort is the typeclass resolution algorithm.
When the elaborator needs an instance of some typeclass, it searches the
instance database. The search is a tournament: each instance is compared
against the current candidate, and the better candidate (more specific
type match) advances. The tournament's winner is the instance the
elaborator selects. The comparison between two candidates is one
unification problem. The tournament's structure ensures that, for $n$
candidate instances, the elaborator performs $O(n)$ unifications, not
$O(n^2)$.

A practical observation: tournament sort is the basis of priority queues
implemented as heaps. A heap is a complete binary tree with the heap
property: each node is at least as large (or as small, depending on the
order) as its children. Insertion and extraction in a heap are both
$O(\log n)$. The heap is the data-structural form of the running
tournament: it always has the current winner at the root, and it can be
updated incrementally as new elements arrive or the current winner is
extracted. Every priority queue, every event-driven simulation, every
Dijkstra's algorithm implementation, every Huffman coding tree, uses a
heap. The heap is the algorithmic generalization of tournament sort.

The chapter's argument is that this generalization is forced by the
chapter's question. To compare records and repeat the comparison reliably,
the algorithmic form must be one of these: sequential scan, binary search,
merge sort, FFT, or tournament sort. The specific form chosen depends on
the question being asked (find one, find all in order, transform between
representations, find the winner without ranking the rest), but the family
of forms is fixed. Other algorithms exist, but they are compositions of
these primitives.

\section{RepeatableProcess in Lean}
\label{se:repeatableprocess-in-lean}

\texttt{Episode2.lean} defines a typeclass \texttt{REPEATABLE} that
formalizes the chapter's central distinction. The class certifies that a
process produces the same output when given the same input, regardless of
when or how often it is called. The definition has the shape:

\begin{leanexcerpt}{REPEATABLE}
class REPEATABLE
    (Value : Type)
    (Carrier : CarrierProcess Value)
    [DISTINGUISHABLE Value Carrier] [ADMISSIBLE Value Carrier]
    [COUNTABLE Value Carrier]       [ENCODED Value Carrier]
    [RESIDUE Value Carrier]
    [b : BINARY Value Carrier]
  where
  repeatable_process : RepeatableProcess Value Carrier
  typical_response   : Trial $\to$ Trial $\to$ Prop := ...
\end{leanexcerpt}

The class carries the full prior cascade through \texttt{BINARY} (each
prior class is a typeclass requirement, not a field). The first field,
\texttt{repeatable\_process}, names the underlying
\texttt{RepeatableProcess} against which the class is calibrated. The
second field, \texttt{typical\_response}, is a binary predicate on
\texttt{Trial}s that records when one trial is an admissible
repetition of another. The proof obligation runs the entire length of
the cascade; consistency under repetition is what
\texttt{typical\_response} certifies.

The typeclass formalizes what every divide-and-conquer algorithm in the
chapter assumes. Binary search assumed that comparing the array element at
position $m$ to the target produces the same Boolean every time. Merge
sort assumed that comparing the front of the left half to the front of
the right half produces the same Boolean every time. FFT assumed that the
butterfly operation produces the same complex outputs every time. Without
\texttt{REPEATABLE}, none of these algorithms can be certified to produce
the correct output.

The proof obligation is stronger than determinism in the runtime sense.
Many programs are deterministic at runtime --- they do not consult random
sources or external state --- but produce outputs that depend on
implementation details like memory layout or floating-point rounding mode.
A \texttt{REPEATABLE} process must produce outputs that are consistent
across implementations: the comparison's value is determined by the
mathematical specification, not by the implementation.

This is why floating-point arithmetic is not trivially
\texttt{REPEATABLE}. The same floating-point computation on two different
CPUs may produce slightly different results due to differences in
extended-precision intermediate registers, fused multiply-add instructions,
or compiler optimizations. The IEEE 754 standard partially addresses this
by specifying exact rounding behavior, but in practice the situation is
murky: some compilers do not strictly follow IEEE 754, and some hardware
provides higher-precision intermediates that can drift the result.

To get a \texttt{REPEATABLE} instance for floating-point comparisons,
either the implementation must be tightly specified (which Lean's
\texttt{Float} type aspires to), or the comparison must be made discrete
(rounded to a fixed precision before comparing). The chapter's algorithms
typically use integer or rational comparisons to sidestep this issue. When
floating-point is required, the \texttt{REPEATABLE} instance comes with
caveats about the implementation.

The companion typeclass \texttt{ComputationalProcess} captures the next
layer: a process that not only is repeatable but produces a numeric
output. \texttt{NUMERIC} is the certificate. Together with
\texttt{REPEATABLE}, it certifies that the output is both deterministic
and admits arithmetic. The chapter's algorithms --- binary search returning
an index, merge sort returning a permutation, FFT returning a complex
vector, tournament sort returning the winner --- all instantiate this
combination. Each produces a numeric output (an integer, a permutation
encoded as integers, a complex vector, a runner's ID) that can be compared
arithmetically with other outputs.

The typeclass \texttt{Study.le} captures the partial order on studies:
one study refines another when its conclusion is at least as specific. A
study that produces a sorted permutation refines a study that produces
only the winner. A study that produces the frequency spectrum of a signal
refines a study that produces only the peak frequency. The partial order
on studies is the partial order on the information content of their
conclusions.

These typeclasses compose to express the chapter's answer. Comparison is
performed by an \texttt{ObservationProcess} that produces \texttt{BINARY}
outputs. Repeated application of the process is admissible only when the
\texttt{REPEATABLE} instance certifies the process. A
\texttt{ComputationalProcess} layers numerical structure: the output is
a number, not just a Boolean. A \texttt{Study} composes multiple trials
into a structured inquiry. The composition produces a permutation, an
index, a spectrum, a ranking --- the algorithmic measurement.

The compiler's enforcement is at elaboration time. If a user writes a
sort function that depends on an unspecified comparison, the
\texttt{REPEATABLE} instance cannot be derived, and the function cannot
be certified. The compiler will not let the user assert that the sort's
output is deterministic without providing the underlying comparison's
determinism proof. This is a non-trivial constraint: many real-world
sorting libraries cannot meet it because they use comparison functions
that consult external state. Those libraries can run, but they cannot be
proven correct in this framework.

The chapter's argument closes with this observation. Divide and conquer
is the algorithmic form of repeatable comparison. The form is forced by
the requirement that the comparison be \texttt{REPEATABLE} (independent
of context) and \texttt{NUMERIC} (producing a structured output). The
specific divide-and-conquer algorithm (binary search, merge sort, FFT,
tournament sort) depends on the question being asked. But the family of
algorithms is the family forced by the chapter's question. Any
comparison that does not factor through this family is either
deterministic-but-not-structurable (a heuristic) or
structurable-but-not-deterministic (an oracle). Both fall outside the
admissibility constraint.

\begin{bridge}

The chapter's question was how an instrument can compare records and
repeat the result without pretending the comparison is the thing
measured.

The answer is divide-and-conquer: the algorithmic family that performs
many small comparisons and composes them into a single structured
output. Binary search finds an index in a sorted array. Merge sort
produces a permutation that orders an input list. FFT transforms a
signal between time and frequency representations. Tournament sort
identifies a winner from a field of contestants. In each case, the
comparisons are the procedure; the structured output is the
measurement. The typeclass \texttt{REPEATABLE} formalizes the constraint
that the comparison produces the same Boolean output on the same input
every time. Without \texttt{REPEATABLE}, the algorithm cannot be
certified.

What remains is the gap between recorded events. Between two
comparisons, the instrument carries some state --- the current
partition, the current subtree, the current frequency band --- but
this state is not yet a record in the chapter's sense. The state is
working memory. To talk about it as a record, the next chapter must
ask what can be said about the state of the algorithm between two
adjacent comparisons, and how the compiler distinguishes admissible
inferences about that state from inadmissible ones. Chapter 3 takes
up the interpolation problem: what the compiler may carry as a
metavariable without resolving it prematurely.

\end{bridge}

\begin{coda}{The Postal Sorting Machine}

The mail-processing center is a long building on the outskirts of the
city. Inside, two parallel conveyor systems carry envelopes and packages
past stations of optical scanners, mechanical diverters, and human
inspectors. The building processes one million pieces of mail per day.
Each piece enters through one of three loading docks and exits through
one of one hundred forty delivery routes. The path each piece takes
through the building is determined by a sequence of comparisons.

The first station is the orientation scanner. An envelope arrives in
arbitrary orientation: front-up, back-up, sideways, upside-down. The
scanner detects which corner has the postage and the address block,
and rotational mechanics in the conveyor turn the envelope so all
envelopes proceed to the next station in the same orientation. This is
not yet a comparison in the chapter's sense; it is preprocessing,
analogous to building the index before sorting.

The second station is the zip-code reader. The optical character
recognition module reads the five-digit zip code at the top of the
delivery address. The reading is a series of small comparisons:
compare the pixel pattern against a template for each digit. The
output is a five-digit number, which becomes the envelope's primary
sort key. The reader achieves about 99.5 percent accuracy. Envelopes
whose zip codes cannot be confidently read are diverted to a manual
station where a human operator types in the zip code from the
displayed image. The diverter is the exception path; the manual
station is the halting-boundary resolution.

After the zip-code reader, the conveyor splits into ten branches. Each
branch handles zip codes starting with one of the ten digits. This is
the first level of divide-and-conquer. The decision at the split is a
single Boolean comparison: is the first digit of the zip code less
than five? If yes, take the upper branch; if no, take the lower branch.
Two further splits at the second level, and so on. By the end of the
five-level split, each branch handles a specific three-digit prefix
range. There are about a thousand such branches.

This is binary search physically realized. The envelope is the target.
The branch system is the search tree. Each split is a comparison. The
total number of comparisons per envelope is $\log_2$ of the number of
final branches, which for a thousand branches is ten. The branch
system processes envelopes in parallel: at any moment, thousands of
envelopes are at different stages of the descent, each one having
made some prefix of the ten comparisons.

The third level of sorting happens within a branch. Each branch's final
chute leads to a local sorting machine that handles four-digit and
five-digit zip code resolution within the branch's three-digit prefix
range. The local machine uses a tournament structure: the envelopes
are ranked against each other by zip code, and the local machine
delivers a sorted stream to the next stage. The tournament is run
continuously: as new envelopes arrive, they enter the bracket and the
local machine re-runs the affected paths. The local machine's output
is a sequence of envelopes in sorted order, one per zip code, ready to
be routed to the appropriate delivery sequence.

The fourth level is the delivery sequence. A delivery sequence is the
order in which a letter carrier visits addresses on a route: down one
side of the street, up the other, into each building in the order the
carrier reaches them. The delivery sequence is a permutation of the
addresses on the route. Each piece of mail's position in the sequence
is determined by another sort, this one keyed by the street number and
apartment number. The fourth-level sort is the most fine-grained: it
operates on individual delivery points, of which there are typically a
few thousand per route.

The fourth-level sort uses merge sort. Multiple input streams (the
envelopes arriving from the local sorting machines) are merged into
a single sorted output, in the delivery sequence order. The merge is
the canonical algorithm: compare the front of each stream, emit the
smallest, advance the corresponding stream's pointer. The merge is
stable: two envelopes for the same address are emitted in their
arrival order, which becomes the order in which the carrier delivers
them (a courtesy to the recipient who may have multiple envelopes
that day).

Throughout all four levels, the comparisons are repeatable. Sending
the same envelope through the machine twice would route it to the
same delivery sequence both times. The comparisons depend only on the
zip code, street number, and apartment number printed on the
envelope, not on the time of day, the operator on shift, or the
weather. The machine's output is deterministic in the input: a
fundamental property of \texttt{REPEATABLE} \texttt{ComputationalProcess}.

The comparison is not the measurement. The measurement is the
delivery sequence: a list of envelopes in the order the carrier will
deliver them. The comparisons are how the machine produces the
delivery sequence; the delivery sequence is what the machine produces.
A million envelopes pass through the machine each day; ten million
comparisons are performed; one hundred forty delivery sequences are
produced. The ratio of comparisons to outputs is roughly $10^7 / 140 =
7 \times 10^4$, which means each delivery sequence is the result of
about seventy thousand comparisons. None of those comparisons is the
sequence itself.

The morning shift manager has been at the center for seventeen years.
She watches the flow on a real-time display: arrivals at the loading
docks, throughput at each station, exceptions diverted to the manual
station, completion times for each delivery sequence. The display
shows aggregate statistics: average sort time, exception rate, branch
utilization. The aggregates are produced from the underlying
comparisons by yet another algorithm: a streaming summarization that
maintains running counts and averages. The aggregates are
measurements, too; the comparisons that produce them are not.

When a problem arises --- a stalled conveyor, a misrouted batch, a
failed optical scanner --- the manager investigates by tracing the
sequence of comparisons backward. She inspects the misrouted
envelopes, reads their zip codes, and checks which comparison led to
the wrong branch. The trace is the audit log of the comparisons: each
envelope's path through the machine is recorded, every decision is
timestamped, every diversion is logged. The audit log is the
\texttt{Study}: the structured record of which trials happened in
which order. The manager uses the study to diagnose the failure.

Today's failure is a misrouted batch of about three hundred envelopes
that should have gone to delivery sequence forty-seven but instead
went to delivery sequence forty-eight. The audit log shows the
problem: at the third-level split for branch forty-seven, the
comparison ``is the fourth digit less than five?'' was producing the
wrong answer for envelopes with fourth digit exactly five. The
optical scanner was misreading the digit five as a six. The
\texttt{REPEATABLE} property had been violated: the comparison was
producing different outputs for envelopes that should have produced
the same output. The fault was at the scanner level, not at the
algorithm level. The algorithm was correct; the input was wrong.

The manager dispatches a technician to recalibrate the scanner. The
recalibration is the next chapter's territory: when the comparison
itself is the source of error, the response is to recalibrate the
instrument. The chapter's algorithms assume \texttt{REPEATABLE}
comparisons; when the assumption fails, the system must intervene to
restore the property. The intervention is not part of the divide-and-
conquer algorithm; it is part of the metrology infrastructure that
supports it.

By mid-afternoon, the technician has replaced the failed scanner
module, recalibrated against a test target, and certified that the
new module produces \texttt{REPEATABLE} comparisons. The misrouted
batch is retrieved, re-scanned, and re-routed through the now-correct
sorting tree. The audit log records both the original misrouting and
the corrective rerouting. The system's record carries both events:
the failure is not erased; it is appended.

At ten o'clock at night, the last delivery sequence is finalized.
One million envelopes have been processed. Ten million comparisons
have been made. One hundred forty delivery sequences have been
produced. The throughput is normal. The exception rate is slightly
elevated due to the morning's scanner failure, but within tolerance.
The shift manager closes the day's audit log and signs off.

The mail-processing center is the most physical possible
embodiment of divide-and-conquer. Every envelope is a record. Every
zip code is a sort key. Every split is a comparison. Every branch is
a recursion level. Every local sort is a tournament. Every merge into
delivery sequences is a stable merge sort. The comparisons are not
the deliveries; the deliveries are what the comparisons produce.

\end{coda}
